% Part: computability
% Chapter: computability-theory
% Section: ce-sets

\documentclass[../../../include/open-logic-section]{subfiles}

\begin{document}

\olfileid{cmp}{thy}{ces}
\olsection{સંગણકીય રીતે પરિગણનીય ગણો}

\begin{defn}
કોઈ ગણ ખાલી હોય અથવા કોઈ સંગણનીય વિધેયનો મૂલ્યગણ હોય તો તે
\emph{સંગણકીય રીતે પરિગણનીય} કહેવાય.
\end{defn}

\begin{history}
સંગણકીય રીતે પરિગણનીય ગણોને \emph{પુનરાવર્તી રીતે પરિગણનીય}
પણ કહે છે. આ મૂળ પરિભાષા છે. આજે બંને નામો તથા ``c.e.'' અને
``r.e.'' સંક્ષેપો સામાન્ય રીતે વપરાય છે.
\end{history}

\begin{explain}
આ વ્યાખ્યાનો અર્થ શું થાય છે અને આ પરિભાષા કેમ યોગ્ય છે તે
વિચારો. વિચાર એવો છે કે $S$ સંગણનીય વિધેય~$f$નો મૂલ્યગણ હોય, તો
\[
S = \{ f(0), f(1), f(2), \dots \},
\]
અને તેથી $f$, $S$ના તત્ત્વોની ``પરિગણના'' કરતું વિધેય છે.
નોંધો કે વ્યાખ્યા મુજબ $f$ વધતું વિધેય હોવું જરૂરી નથી; એટલે
પરિગણના ચડતા ક્રમમાં હોવી જરૂરી નથી. હકીકતમાં, $f$ એક-એક
હોવું પણ જરૂરી નથી; એટલે $S$ની પરિગણનામાં $f(0)$, $f(1)$,
$f(2)$, \dots{} નું પુનરાવર્તન ચાલે છે. દાખલા તરીકે,
અચળ વિધેય $f(x) = 0$, ગણ $\{ 0 \}$ની પરિગણના કરે છે.
\end{explain}

દરેક સંગણનીય ગણ સંગણકીય રીતે પરિગણનીય છે. આ જોવા માટે
માનો કે $S$ સંગણનીય છે. $S$ ખાલી હોય, તો વ્યાખ્યા મુજબ તે
સંગણકીય રીતે પરિગણનીય છે. અન્યથા $S$નું કોઈ તત્ત્વ $a$ લો.
$f$ને નીચે પ્રમાણે વ્યાખ્યાયિત કરો:
\[
f(x) =
\begin{cases}
x & \text{જો $\Char{S}(x) = 1$} \\
a & \text{અન્યથા.}
\end{cases}
\]
ત્યારે $f$ સંગણનીય વિધેય છે અને $S$ એ $f$નો મૂલ્યગણ છે.

\end{document}
